\begin{statement}{B.14b}{Proved}{The unscaled exceptional-prime witness has an integral lift through class fourteen}
Let $F=\langle x,y\rangle$, $N_{14}=F/\gamma_{15}F$,
and $z=(xy)^{-1}$. There is an integral Hall group word
$f\in\gamma_7N_{14}$ with leading term the specified $q_7$ of
B.12, without scaling, such that
\[
 f(y,x)f(x,y)=1,\qquad
 f(z,x)f(y,z)f(x,y)=1
 \quad\text{in }N_{14}.
\]
The complete affine system in all Hall exponents of weights
eight through fourteen has $5030$ Lyndon equations and $2497$
unknowns, exact rational rank $2066$, and nullity $431$.
Its integral solutions form a nonempty affine lattice of rank
$431$. In particular the unscaled finite jet exists over
$\mathbb Z_p$ for every prime $p$, including $2$ and $3$.
This assertion does not establish an infinite lift.
\end{statement}
\begin{proof}
Use the commutators $C_w$ and Magnus conventions of B.14a,
ordering first by weight and then lexicographically. Write
\[
 f=\left(\prod_{|w|=7}C_w^{Q_w}\right)
       \left(\prod_{8\leq |w|\leq14}C_w^{u_w}\right),
\]
where the $Q_w$ are the fixed coefficients of $q_7$.
Every one of the $2497$ variables $u_w$ is included; no earlier
free parameter is specialized to the previously chosen
class-thirteen certificate.

The quadratic terms in weight fourteen must now be retained.
For any of the four cycle substitutions, put $a_w=C_w-1$
and denote its degree-seven part by $a_{w,7}$. The fixed
weight-seven Hall word has deviation $S+R$, where
\[
 S=\sum_{|w|=7}Q_wa_w,
 \qquad
 R=\sum_w\binom{Q_w}{2}a_{w,7}^2
       +\sum_{v<w}Q_vQ_wa_{v,7}a_{w,7}.
\]
These formulas hold for negative as well as positive integral
exponents, and $R$ is homogeneous of degree fourteen. For
$|w|\geq8$, however,
$C_w^{u_w}-1=u_w(C_w-1)$ through degree fourteen.
An interaction of an unknown factor with a fixed weight-seven
factor starts in degree $15$, and an interaction of two
unknown factors starts in degree $16$. Thus the complete
system remains \emph{affine linear} in all the $u_w$.

Index the substitutions by
$(x,y),(y,x),(z,x),(y,z)$, numbered $0,1,2,3$.
If $F_i$ is the fixed Hall-word deviation and $L_i$ its
degree-seven part, the two fixed cycle defects are exactly
\[
 F_1+F_0+L_1L_0,\qquad
 F_2+F_3+F_0+L_2L_3+L_2L_0+L_3L_0.
\]
The code includes every one of these ordered products.
The resulting quadratic corrections have respectively
$2048$ and $9200$ nonzero associative-word coefficients.
The complete system $Au=b$ is stored in
\texttt{code/g\_exceptional\_matrix\_14.json.gz}; its
entire part through weight thirteen agrees exactly with the
independent matrix of B.14a.

Selecting Lyndon-word rows loses no equation. Each cycle
product is group-like; formal binomial powers over
$\mathbb Q$ give the same formulas for rational kernels.
Its first nonzero homogeneous
deviation, if present, would therefore be an ordinary
homogeneous free Lie polynomial over $\mathbb Q$.
The triangular Lyndon basis detects such a polynomial on
Lyndon words. Induction on the first nonzero degree proves
the assertion. The same argument with zero weight-seven
term applies to the homogeneous matrix kernel, so its
rank is also the rank of the complete associative-word
system.

For the exact rank and integral solve, first eliminate
$1776$ variables using only pivots $+1$ and $-1$.
Each eliminated equation is recorded as
$u_j=b_i-\sum_{k\ne j}a_{ik}u_k$.
These operations are reversible over $\mathbb Z$ and
preserve every integral choice of the remaining variables.
The resulting $2787\times721$ matrix has exact rank $290$;
a modular row basis modulo $1000003$ has the same rank.
Consequently the full rank is $1776+290=2066$.
The exact residual system and the back-substitution equations
are stored in
\texttt{g\_exceptional\_reduction\_14.json.gz}.

We establish integral consistency by an exact finite local
calculation, without assuming that a rational solution is
integral. Select $290$ independent residual rows and write
their system as
\[
 Bv+Ct=y,\qquad
 v=\beta-St,\quad
 \beta=B^{-1}y,\quad S=B^{-1}C,
\]
where $B$ is an invertible rational minor and the $431$
coordinates $t$ remain free. For a prime $p$, let $d$ be
the largest $p$-denominator exponent among the entries of
$\beta$ and $S$. For $t\in\mathbb Z_p^{431}$, the vector
$v$ is $p$-integral precisely when
\[
 p^dSt\equiv p^d\beta\pmod{p^d}.
\]
The coefficients belong to $\mathbb Z_{(p)}$ and are
reduced by inverting their prime-to-$p$ denominators.
This is an exact finite congruence criterion: an integral
representative of any solution for $t$ gives a rational
solution with $p$-unit denominator.

The finite congruence solver in
\texttt{g\_exceptional\_local14.py} inverts only unit
pivots. When every remaining coefficient is divisible by
$p$, it checks the same divisibility on every right side
before dividing the whole residual system and its modulus
by $p$. Back-substitution uses the modulus of each recorded
unit equation. The resulting vector is checked against
the original finite congruences and then against
\emph{every} residual rational equation.

For both $p=2$ and $p=3$, the selected minor gives $d=1$.
The two exact local solutions are $U_2/d_2$ and $U_3/d_3$,
where the $U_p$ are integral and $p\nmid d_p$.
The denominators have respectively $129$ and $108$ decimal
digits and are coprime. The certificate records integers
$s_2,s_3$ with $s_2d_2+s_3d_3=1$.
Thus $s_2U_2+s_3U_3$ solves the residual system integrally.
Integer back-substitution recovers all $2497$ coordinates,
and every row of the original $Au=b$ is checked exactly.
This combines affine correction vectors while keeping
$q_7$ fixed; it does not scale the quadratic right side or
assert that scaling it scales the leading group term.

The resulting group word has $2079$ nonzero Hall
coordinates, including the seven fixed weight-seven
coordinates. Its largest exponent has $242$ decimal
digits; all integers, local minors, congruence solutions,
and B\'ezout coefficients are retained in
\texttt{g\_exceptional\_local\_14.json}.
Finally \texttt{g\_exceptional\_14.py} reconstructs all
four substituted Hall group words and checks every
associative-word coefficient of the two group products
through degree fourteen. Both defects are exactly zero.
The fixed quadratic terms are also checked independently
by repeated multiplication and inverse factors.
As in B.14a, a nonzero lower-central component through
degree fourteen would give a nonzero leading Magnus term.
The verified identities therefore hold in $N_{14}$.

Since an integral solution exists, all integral solutions
are its translate by $\ker(A:\mathbb Z^{2497}\to
\mathbb Z^{5030})$, a free lattice of rank $431$.
Passing the displayed word to the pro-$p$ nilpotent quotient
preserves both identities. No compatibility at uncomputed
higher classes is inferred.
\end{proof}
